\section{Introduction}

Reversibility turns periodic-orbit equations into intersections of fixed
curves of involutions.  This viewpoint is classical and remains useful in
modern studies of symmetric periodic points
\citep{Kang2014,YamaguchiTanikawa2005}.  It also creates a precise
effectivity question: when does an algebraic symmetry-line closure count
reduced primitive intersections?

Two predecessor results make that question nontrivial here.  First, a
certified four-state survivor of $H_6$ is mixing and uniformly hyperbolic.
Its primitive reversible cycles have exact entropy $\frac12\log\varphi$.
Second, the odd mixed-axis closure satisfies an exact divisor law and has
formal primitive degree of entropy $\frac12\log2$.  Finite quotients are
reduced through period fifteen, but general projective dynatomic effectivity
\citep{Hutz2010} does not automatically attach to this birational
symmetry-line slice.

We separate the physical and ambient questions.  Our main theorem closes the
physical one at every odd period.

\begin{theorem}[Physical reflection transversality]\label{thm:main}
Every primitive odd mixed-axis closure root belonging to the certified H6
survivor is simple.  The number $P_n$ of such roots is
\[
P_n=\sum_{d\mid n}\mu(n/d)F_{(d+3)/2},
\]
and every one has coefficient $+1$ in the formal mixed-axis dynatomic
divisor.
\end{theorem}

Consequently the physical roots are exponentially sparse in formal degree.
This is simultaneously a positive local-effectivity theorem and a warning:
hyperbolicity of one real survivor cannot certify the exponentially larger
ambient algebraic population.
